\chapter{Computer vision approaches for detection of the abnormal fenotype in yeast cells}

Microscopy enables capturing the image data from cell populations. 

\subsection{Image analysis driven single-cell analytics for systems microbiology}

~\cite{balomenos2017image}

\subsection{Nanoparticle Recognition on Scanning Probe Microscopy Images Using Computer Vision and Deep Learning}

~\cite{okunev2020nanoparticle}

\subsection{Deep learning microscopy}

~\cite{rivenson2017deep}

\subsection{Machine learning and computer vision approaches for phenotypic profiling}

~\cite{grys2017machine}

\subsection{Iterative thresholding based image segmentation using 2D improved Otsu algorithm}

~\cite{devi2015iterative}

\subsection{A quantitative analysis of cell bridging kinetics on a scaffold using computer vision algorithms}

~\cite{lanaro2021quantitative}

\subsection{Computer vision for high content screening}

~\cite{kraus2016computer}

\subsection{YeastNet: Deep-Learning-Enabled Accurate Segmentation of Budding Yeast Cells in Bright-Field Microscopy}

Deep-learning approach for accurate segmentation of the yeast cell bright-field microscopy data~\cite{salem2021yeastnet}. The convolutional network is based on the U-Net semantic segmentation architecture~\cite{ronneberger2015u}.

\subsection{U-Net: Convolutional Networks for Biomedical Image Segmentation}

~\cite{ronneberger2015u}